\chapter{Grandes volúmenes: tres cajas y Spark}
\label{cap:s-cajas}
\textit{Materia: Grandes volúmenes de datos. Pregunta: ¿cómo se guardan, se limpian y se juntan millones de datos de
fuentes que no se parecen entre sí?}

\section{El problema}
Las fuentes llegan en formatos distintos: archivos Excel con notas al pie (SECTUR), archivos de texto con un formato
propio (las tormentas de la NOAA), respuestas de una página web (SITUR-Q), CSV comprimidos con acentos raros (el DENUE).
Algunas son historia (el Censo de 2020) y otras cambian cada semana (la ocupación). Hay que juntarlas sin perder nada y sin
mezclar lo crudo con lo limpio.

\section{Las tres cajas}
\begin{intuicion}
Es como el agua que pasa por un filtro de tres etapas. La primera caja guarda el agua tal como llegó del río; la segunda,
el agua filtrada; la tercera, el agua lista para servir. Si algo sale mal en el filtro, siempre se puede volver al agua
original.
\end{intuicion}

\begin{enumerate}
\item \textbf{Caja de bronce} (\codigo{datos/bronze/}): los 353 archivos tal como llegaron. \textbf{Nunca se tocan.}
\item \textbf{Caja de plata} (\codigo{datos/silver/}): los datos limpios, con nombres de columna en español, tipos
  correctos y marcas de calidad. Hay un archivo de limpieza por fuente (12), porque cada una tiene sus propias trampas.
\item \textbf{Caja de oro} (\codigo{datos/gold/}): lo que usan los modelos y la página: la tabla central de hechos y los
  resultados de cada fase.
\end{enumerate}

\begin{porque}
Si se limpiara al descargar, el original se perdería, y si la regla de limpieza estaba mal, no habría forma de rehacerla.
Con tres cajas, cualquier limpieza se puede repetir desde cero. Esta forma de organizar se llama \emph{arquitectura
medallón} y es la que usan las empresas que manejan muchos datos.
\end{porque}

\section{La huella digital de cada archivo}
Cada archivo de la caja de bronce tiene una \textbf{huella} (SHA-256): un código de 64 letras y números que se calcula con
su contenido. Si alguien cambia una sola coma del archivo, la huella cambia por completo. Una prueba automática recalcula
las 353 huellas y avisa si alguna no coincide. Así se demuestra que los datos crudos son los mismos que se descargaron.

Todo queda anotado en un \textbf{manifiesto} (\codigo{datos/bronze/MANIFIESTO.csv}): el nombre del archivo, su huella, su
tamaño, de qué dirección vino, cuándo se bajó y cuántos registros tiene.

\section{Spark: muchas manos al mismo tiempo}
\begin{intuicion}
Revisar 6.1 millones de negocios es como calificar 6.1 millones de exámenes. Una persona sola tardaría muchísimo. Doce
personas, cada una con un montón, terminan mucho antes. Spark hace eso: parte la tabla en montones (\emph{particiones}) y
pone a trabajar a todos los núcleos de la computadora a la vez.
\end{intuicion}

Cinco ideas de Spark que conviene poder explicar:
\begin{enumerate}
\item \textbf{Un jefe y varios trabajadores.} Un programa planea el trabajo y los trabajadores lo hacen. Aquí todo corre
  en una sola computadora, con sus 12 núcleos. Es el mismo programa que correría en un centro de datos con cientos de
  computadoras: solo cambia la dirección del jefe.
\item \textbf{Tablas repartidas.} Cada pedazo de la tabla lo procesa un núcleo distinto.
\item \textbf{Primero el plan, luego el trabajo.} Spark no ejecuta cada instrucción en cuanto se escribe: arma el plan
  completo y lo optimiza. Por ejemplo, si al final solo se necesitan 10 columnas, no lee las otras.
\item \textbf{Juntar lo repartido cuesta.} Quitar negocios repetidos obliga a juntar en el mismo núcleo todos los renglones
  con la misma clave. Es la parte más cara y por eso se hace una sola vez.
\item \textbf{Parquet por carpetas.} Spark guarda cada tabla en formato Parquet, que se lee por columnas y viene
  comprimido, y la reparte en carpetas (por ejemplo, una por estado). Para leer solo Quintana Roo no hace falta abrir los
  otros 31 estados.
\end{enumerate}

\begin{ejemplo}[ · el DENUE completo]
Spark leyó los 33 archivos del DENUE (el Estado de México viene en dos partes) en cerca de un minuto y medio y:
\begin{itemize}
\item contó 6,138,075 negocios, con 0 repetidos y 0 coordenadas fuera de México;
\item quitó la razón social, el teléfono y el correo;
\item marcó cuáles son turísticos según su giro oficial (hoteles, restaurantes, agencias de viajes, transporte
  turístico, museos y esparcimiento): 875,667 en el país, \textbf{13,663 en Quintana Roo}.
\end{itemize}
\end{ejemplo}

\begin{teprofe}
\emph{``¿Para qué Spark si cabe en tu computadora?''} Porque el incidente de Big Data pide una arquitectura que siga
funcionando cuando los datos \emph{no} quepan. El código de Spark es el mismo para 6 millones que para 600 millones de
registros: solo se agregan computadoras. Además, el DENUE nacional (574 MB comprimido) ya es pesado para leerlo de golpe con
herramientas comunes.
\end{teprofe}

\section{Reglas de limpieza que salieron de los datos}
Cada regla se comprobó con números antes de aplicarla:
\begin{itemize}
\item \textbf{Tren Maya}: algunos destinos traen dos renglones por mes (uno por estación) y \textbf{se suman}. Se comprobó
  que la suma cuadra exacto con el total que publica el gobierno. En enero de 2025: Chetumal ($24+3{,}478$) + Bacalar
  ($213+3{,}369$) = 7,084, igual que el total de la zona.
\item \textbf{Ceros imposibles = hueco} (capítulo anterior).
\item \textbf{Una versión por periodo}: SECTUR corrige sus cifras preliminares, y un mismo periodo aparece en varios
  archivos. Gana el más reciente. Hubo 2,804 cifras semanales corregidas.
\item \textbf{Repetidos}: el INAH traía repetido el bloque de extranjeros de septiembre de 2025, una copia con cifras y
  otra en ceros. Se conserva la de cifras. Si algún día aparecen dos cifras distintas para lo mismo, el programa se
  detiene en lugar de escoger una.
\item \textbf{Feriados sin tipo de cambio}: se quedan vacíos. Copiar el día anterior inventaría una cotización que no
  existió.
\end{itemize}

\section{El almacén: una biblioteca con fichero}
Al final, todo se puede consultar en \textbf{DuckDB}, una base de datos que vive en un solo archivo
(\codigo{datos/gold/torre.duckdb}). Lee las tablas de las cajas de plata y oro sin copiarlas y responde en milisegundos,
sin servidor y sin internet.

La tabla central (\emph{hechos}) tiene un renglón por lugar, mes y variable, con su fuente: 3,287 renglones de 15 lugares
y 192 meses. \textbf{Un hueco no tiene renglón}; nunca se rellena con cero. A los lados están las tablas de lugares y de
fechas. Esta forma se llama \emph{modelo estrella}: la tabla de hechos en el centro y las descripciones alrededor.

Además se genera solo un \textbf{diccionario de datos} de 56 tablas: cada columna con su tipo, su porcentaje de vacíos, un
ejemplo real y su significado. Si una columna no tiene significado escrito, el programa se detiene.

\section{Datos que llegan cada semana: la Torre con Spark Streaming}
El incidente de Big Data pide una arquitectura de \textbf{baja latencia}: que la campaña se ajuste \emph{mientras} corre.
Para eso, la Torre usa \textbf{Spark Structured Streaming}, la parte de Spark que procesa datos conforme van llegando.

\begin{intuicion}
Es como ver una película cuadro por cuadro. Cada semana real (de enero de 2022 a julio de 2026) se guarda como un archivo,
y Spark los lee de uno en uno, en orden, como si fueran llegando. En cada uno decide si pausa un anuncio.
\end{intuicion}
Se reprodujeron \textbf{239 semanas en 239 lotes}, en orden. Spark anota en una bitácora (\emph{checkpoint}) cuál fue el
último archivo leído: si el proceso se cae, sigue donde se quedó sin repetir ni saltarse semanas. Y si mañana llegara una
fuente semanal real del sur, se cambia la carpeta y las reglas siguen igual.
